% ============================================================
%  第三部分：判定质量的验证体系
% ============================================================

\section{判定质量的验证体系}
\label{sec:validation}

这是最容易被跳过的一章，也是区分“研究玩具”与“生产系统”的分水岭。Regime 建模有一个结构性困难：\textbf{没有真实标签}。我们从未“看见”过市场的真实状态——所有标签都来自模型或规则本身。一旦评估标准退化为“回测赚钱就行”，系统就同时失去了诊断能力和纠错能力：赚了不知道为什么，亏了也不知道为什么。

本章建立三层间接证据体系（模型内证据、结构证据、经济证据），给出前瞻偏差的完整检查清单与稳定性检验方案，最后汇总为可执行的评估工作流。

\subsection{无标签世界的三类间接证据}

\begin{table}[htbp]
\centering
\footnotesize
\begin{tabular}{p{2.6cm}p{4.6cm}p{5.4cm}}
\toprule
\textbf{证据类型} & \textbf{回答的问题} & \textbf{典型工具} \\
\midrule
模型内证据 & 状态“分得开”吗？画像清晰吗？ & 状态画像表、分离度、持续期、转移矩阵 \\
结构证据 & 结论“经得起扰动”吗？ & 子样本重估、种子扰动、块 bootstrap、ARI \\
经济证据 & 对决策“有用”吗？ & 条件化策略样本外表现、成本、换手 \\
\bottomrule
\end{tabular}
\caption{无标签条件下的三类间接证据}
\label{tab:validation-evidence}
\end{table}

三类证据缺一不可：只有经济证据 $=$ 可能靠偶然与泄漏赚钱；只有模型内证据 $=$ 可能是漂亮的拟合但无经济含义；只有结构证据 $=$ 稳定但无用。完备的评估必须三者交叉。

\subsection{第一层：内部诊断清单}

对训练好的任何状态模型（HMM、JM、规则法），逐项检查：

\begin{enumerate}[itemsep=3pt]
\item \textbf{状态画像表}。为每个状态列出：样本占比、均值收益、波动率、平均持续期、代表性时段。判据：画像之间应\textbf{有清晰的语义差异}（高波 vs 低波、正收益 vs 负收益），而不是“三个状态长得差不多”。若两个状态画像重叠，直接减 $K$。
\item \textbf{分离度}。两两状态的分布距离（如 Bhattacharyya 距离、或简单地看收益率/波动率的分离倍数）。分离度低的边界区域会出现高频抖动式切换。
\item \textbf{持续期分布}。与先验期望对照：过短（如平均 < 3 期）说明模型在噪声上抖动；过长或单状态占比 $> 90\%$ 说明状态坍缩（有效 $K$ 小于名义 $K$）。
\item \textbf{转移矩阵结构}。对角优势（持久性）应当是常态（$a_{ii}$ 显著大于 $1/K$）；若对角无优势，模型没有学到“状态”这一结构。
\item \textbf{多起点似然分布}。多个随机初始化下对数似然的最大值与中位数的差距：差距大说明局部最优问题严重，必须保留多起点与择优机制。
\item \textbf{似然—状态数曲线}。$K$ 从 1 增至 6，观察似然改善的“肘部”位置；结合第~\ref{sec:essence}~节的告诫：$K$ 是决策分辨率选择，不宜超过数据能支撑的分离度。
\end{enumerate}

\subsection{第二层：前瞻偏差——四种来源与完整检查清单}

前瞻偏差（look-ahead bias）在 Regime 系统中比在普通策略中更隐蔽，因为它发生在四个位置：

\begin{enumerate}[itemsep=3pt]
\item \textbf{平滑泄漏}。用全样本平滑概率 $\gamma_t$（或全样本 Viterbi 路径）做回测——第二章已证明 $\gamma_t$ 依赖未来数据。表现为：“状态几乎完美地在转折点当天切换”。
\item \textbf{参数泄漏}。用全样本估计的参数 $\hat{\btheta}$（而非滚动/扩展窗口逐步重估）计算历史状态。表现为：早期时段的状态也拥有“后来的知识”。
\item \textbf{特征泄漏}。状态判定所用的特征本身包含未来信息：如全样本 z-score 标准化、全样本分位数阈值、全样本分组的“高波状态”定义。
\item \textbf{执行泄漏}。信号在 $t$ 日收盘形成，却假设 $t$ 日收盘即可按收盘价成交。严谨设定应至少延迟到 $t+1$ 开盘/收盘~\cite{shu2024a}。
\end{enumerate}

\begin{remark}[泄漏审计实验：一个必须做的对照]
最有效的审计方法不是读代码，而是做\emph{对照实验}：把同一套策略跑两遍——一遍用严格的滤波 $+$ 滚动参数（合规版），一遍故意用平滑 $+$ 全样本参数（泄漏版）。两者绩效之差，就是系统对前瞻偏差的敏感度。若泄漏版比合规版好得\emph{不成比例}（例如 Sharpe 翻倍），说明该系统高度依赖未来信息，必须彻底重构推理管线的数据流。这个实验成本极低且证据力极强，建议纳入常规回归测试。
\end{remark}

\subsection{第三层：稳定性检验}

\begin{itemize}[itemsep=3pt]
\item \textbf{子样本重估}。把样本切成 3--5 个时段分别重估，观察状态画像与持续期的漂移。画像维持 $\Rightarrow$ 结构稳定；画像重构 $\Rightarrow$ 历史校准失效（警示结构演化已侵蚀历史）。
\item \textbf{参数扰动}。对初始化种子、窗口长度、特征子集做扰动，用调整后兰德指数（ARI，衡量两次分割的一致性）度量状态序列的稳定程度。ARI 显著低于 $0.7$ 时，不应把该系统的单次输出当作“事实”上报。
\item \textbf{块 bootstrap}。对转移矩阵元、平均持续期、状态均值构造置信区间（块采样保留时序依赖）。报告点估计而不给区间，是研究文档中常见的过度自信。
\item \textbf{身份一致性}。跨滚动窗口的状态匹配率——即第~\ref{sec:frontier}~节 Wasserstein 追踪要解决的问题。即便没有采用前沿方法，也应对每个状态维护“画像指纹”并检查逐窗匹配：匹配率骤降的窗口，是整个系统最需要人工复核的窗口。
\item \textbf{变点不确定性}。对任何“结构断裂”的陈述，报告区间而非精确日期（呼应第~\ref{sec:boundaries}~章之前的讨论与 Casini--Perron~\cite{casini2022}）。
\end{itemize}

\subsection{经济价值评估：口径与陷阱}

到这一层，评估看起来“传统”了，但 Regime 系统有三个特有陷阱：

\begin{enumerate}[itemsep=3pt]
\item \textbf{必须与同风险预算的静态基准对比}。比较对象不是“无策略”，而是“相同波动率目标下的固定姿态”。否则你度量的是杠杆差异而不是状态识别的价值。
\item \textbf{必须计入状态切换的调仓成本}。状态路由的换手远高于静态配置——每一次误报切换都是实打实的磨损（这也是第~\ref{sec:frontier}~节为何“低换手控制”是前沿主线之一）。
\item \textbf{必须处理多重检验}。研究过程中尝试的每一个窗口长度、状态数、阈值，都是一次检验；从成百次试验中选出“最好”的结构，其绩效的统计显著性被系统性高估。Harvey et al.~\cite{harvey2016} 给出了这类“数据窥探”在因子研究中的量化框架（$t$ 值门槛需远高于 $2$），同一逻辑原样迁移到 Regime 系统中：报告的绩效应在“全部试验集合”的意义上打折扣。
\end{enumerate}

\subsection{汇总：可执行的评估工作流}

\begin{enumerate}[itemsep=3pt]
\item \textbf{内部诊断}（第~\ref{sec:validation}~节第一层）：画像表、分离度、持续期、多起点似然。不通过则回到建模，不进入回测。
\item \textbf{泄漏审计}：合规版 vs 泄漏版对照实验；检查清单逐项签字（每个 $t$ 的状态与参数只依赖 $Y_{1:t}$；特征与阈值滚动生成；信号—执行至少间隔 1 期；成本按实际换手计）。
\item \textbf{稳定性检验}：子样本重估、种子/参数扰动（ARI）、块 bootstrap 区间、身份一致性匹配率。
\item \textbf{经济价值}：同风险预算基准对比、含成本与延迟、多重检验折扣。
\item \textbf{Go / No-Go 判据}：三项证据全部通过的组合进入生产灰度；任何单项失败的系统只允许作为“研究输入”，不得进入决策链路。
\end{enumerate}

\begin{table}[htbp]
\centering
\footnotesize
\begin{tabular}{p{3.6cm}p{5.4cm}p{3.6cm}}
\toprule
\textbf{失败模式} & \textbf{典型症状} & \textbf{对应检验} \\
\midrule
状态抖动 & 平均持续期 < 3 期、切换频繁 & 持续期分布、分离度 \\
状态坍缩 & 单状态占比 > 90\%、似然平坦 & 画像表、转移矩阵 \\
前瞻偏差 & 合规/泄漏对照差异巨大 & 泄漏审计实验 \\
历史失效 & 子样本画像重构、ARI 低 & 子样本重估、身份一致性 \\
绩效幻觉 & 含成本后优势消失 & 成本计入、多重检验折扣 \\
\bottomrule
\end{tabular}
\caption{失败模式与对应检验的映射}
\label{tab:failure-modes}
\end{table}

\keynote{Regime 模型没有真实标签，只有三类间接证据：模型自洽、结构稳定、经济价值。只报告一类的评估——尤其是只报告回测绩效——不是评估，是广告。}